problem:  
Let \( (X,\omega,\Omega) \) be a compact Calabi–Yau threefold, and let \(L_0 \subset X\) be a compact **immersed** special Lagrangian submanifold calibrated by \(\operatorname{Re}(\Omega)\) with a single **transverse double point** \(p \in L_0\).  Suppose the local branches of \(L_0\) at \(p\) intersect at an **angle** \(\theta \in (0, \pi)\) and satisfy the **cohomological balancing condition** required for Joyce’s desingularization theory.  

Joyce constructed small embedded special Lagrangians \(L_\epsilon\) obtained by gluing in a **Lawlor neck** near \(p\); the existence of such desingularizations (for some small \(\epsilon\)) is known under the given local conditions.  

**Problem:**  
In the above setting, is it true that:  

1. There exists a **unique (up to reparametrization)** smooth one-parameter family of embedded special Lagrangians  
   \[
   \{L_\epsilon\}_{\epsilon>0}, \quad \text{with } L_\epsilon \to L_0 \text{ weakly as } \epsilon \to 0,
   \]
   where the convergence is in the sense of integral calibrated varifolds, and the parameter \(\epsilon\) corresponds to the scaled neck size of the gluing construction?  

2. This family exhausts all small desingularizations of \(L_0\) near the given double point—i.e., every embedded special Lagrangian sufficiently \(C^1\)-close to \(L_0\) and representing the same homology class arises from precisely one gluing parameter \(\epsilon>0\)?  

3. Consequently, the local moduli space of (possibly singular) special Lagrangians in the homology class \([L_0]\) admits a **natural manifold-with-boundary structure**
   \[
   \overline{\mathcal{M}}_{\mathrm{SLag}}([L_0]) \;\cong\; [0,\epsilon_0) \times \mathcal{V},
   \]
   where the boundary at \(\epsilon=0\) corresponds to the singular limit \(L_0\) and the interior corresponds to embedded desingularizations \(\{L_\epsilon\}\)?  

Why is it a "good" problem:  
This problem isolates one of the central analytic mysteries in special Lagrangian geometry: whether the space of special Lagrangians can be locally compactified by adding singular (immersed) objects as genuine boundary points, in a manner **analytically controlled** by a geometric gluing parameter.  

It refines Joyce’s desingularization theory into a precise and minimal conjecture about **local moduli topology**—essentially asking whether the expected “one-dimensional smoothing direction” arising from linearized deformation theory actually produces a canonical manifold-with-boundary extension once nonlinear effects and global constraints are accounted for.  

The question demands a synthesis of:  

- nonlinear elliptic analysis for the special Lagrangian equation under singular perturbations;  
- geometric measure theory (calibrated varifold limits); and  
- moduli identification issues comparing analytic gluing families.  

A positive resolution would establish the first analytic model for **real conifold transitions**, bridging geometric analysis, mirror symmetry, and moduli compactification theory. Its simplicity—“does a single-neck desingularization parameter canonically compactify the special Lagrangian moduli?”—contrasts with the technical depth required to resolve it. This balance of conceptual clarity, cross-field connection, and analytic subtlety makes it an excellent and truly “good” research problem.